\section*{Notch-T}

\noindent\colorbox{orange}{
  \begin{minipage}{1.0\linewidth}
    \paragraph{Richiesta (a)}
    Sulla base di semplici argomenti di natura qualitativa (ovvero senza effettuare il calcolo
    della funzione di trasferimento, come richiesto al punto successivo), mostrare che il
    circuito svolga effettivamente la funzione di notch, discutendone la risposta ai limiti del
    dominio delle frequenze reali.
  \end{minipage}
}
\begin{itemize}
\item
  Per $f \to 0$ i condensatori nel circuito ``aprono'' i loro rami. Di conseguenza il circuito
  finale si tratta di un op-amp in configurazione di \textit{inseguitore di tensione}, la tensione
  $V_S$ attraversa $R_1$ e $R_2$ in serie, nell'ipotesi di ``grossa'' impendenza in ingresso
  all'op-amp tali resistenze possono essere trascurate e nel regime lineare di operazione
  dell'op-amp vale
  \begin{equation*}
    V_{OUT} = A_d(V_S - V_{OUT}) \implies \colorbox{blu}{\boxed{V_{OUT} = \frac{A_d}{1+A_d}V_S}}
  \end{equation*}
\item
  Per $f \to \infty$ i condensatori $C_1$ e $C_2$ mettono ``in corto-circuito'' $V_S$ e $V_+$,
  nel regime lineare di operazione dell'op-amp vale
  \begin{equation*}
    V_{OUT} = A_d(V_S - V_{OUT}) \implies \colorbox{blu}{\boxed{V_{OUT} = \frac{A_d}{1+A_d}V_S}}
  \end{equation*}
\end{itemize}
La caratteristica del notch \'e quella dell'eliminazione di una banda all'uscita dal circuito,
data la topologia del circuito la condizione equivalente \'e quella dell'annullamento della
tensione all'ingresso \underline{non}-invertente $V_+$.
\begin{itemize}
\item
  \textit{Aumentando} la frequenza del segnale $f \to 0$, durante la caduta/salita di potenziale
  fra $R_1$ e $R_2$ il condensatore $C_3$ si scarica/carica, immettendo/assorbendo corrente da
  $R_2$ e di conseguenza \textit{ritardando} il segnale.
\item
  \textit{Diminuendo} la frequenza del segnale $f \to \infty$
  la carica del condensatore $C_1$ impiega tempo non nullo durante cui $C_2$ si carica pi\'u
  lentamente e di conseguenza \textit{ritardando} il segnale.
\end{itemize}
\noindent\colorbox{blu}{
  \begin{minipage}{1.0\linewidth}
    Siccome agli estremi del dominio di frequenza il segnale \'e senza attenuazione e per piccole
    variazioni di frequenza esso tenda ad attenuarsi \'e ``ragionevole'' assumere l'esistenza di un
    minimo anche detto \textit{notch}.
  \end{minipage}
}
\paragraph{Richiesta (b)}
Derivare analiticamente la funzione di trasferimento $A(s)=V_{OUT}(s)/V_S$ nell'ipotesi che
l'operazionale sia ideale ed assumendo per semplicit\'a che $R_1=R_2=R=10 kW$, $R3=R/2$,
$C_1=C_2=C=10 nF$ e $C_3=2C$. Si suggerisce di procedere:
\begin{enumerate}
\item[i.]
  considerando lo schema semplificato del circuito riportato in calce;
\item[ii.]
  ottenendo il suo equivalente per effetto della sostituzione di Th\'evenin dei blocchi A
  e B;
\item[iii.]
  quindi determinando la tensione nel nodo $V_+$ in termini di $V_S$ e $V_{OUT}$ ed in funzione
  della variabile complessa adimensionale $u = sRC$;
\item[iv.]
  infine, mettendo in relazione $V_+$ e $V_{OUT}$.
\end{enumerate}
\paragraph{Risposta (b.ii)}
\begin{itemize}
\item[a.]
  \begin{equation*}
    \begin{gathered}
      sC(V_S - V_A) = \frac{V_A - V_{OUT}}{R/2} \\
      sRC(V_S - V_A) = 2(V_A - V_{OUT}) \\
      sRCV_S + 2V_{OUT} = (2 + sRC)V_A \\ \implies
      V_{Th, A} \equiv V_A = \frac{sRCV_S + 2V_{OUT}}{2 + sRC}
    \end{gathered}
  \end{equation*}
  \begin{equation*}
    Z_{th, A} \equiv \frac{R}{2 + sRC}
  \end{equation*}
\item[b.]
  \begin{equation*}
    \begin{gathered}
      \frac{V_S - V_A}{R} = 2sCV_A \\
      V_{Th, B} \equiv V_A = \frac{V_S}{1 + 2sRC}
    \end{gathered}
  \end{equation*}
  \begin{equation*}
    Z_{Th, B} \equiv \frac{R}{1 + 2sRC}
  \end{equation*}
\end{itemize}
\paragraph{Risposta (b.iii)}
\begin{equation*}
  \begin{gathered}
    V_{Th, A} = \frac{uV_S + 2V_{OUT}}{2 + u} \qquad
    Z_{Th, A} = \frac{R}{2 + u} \\
    V_{Th, B} = \frac{V_S}{2 + u} \qquad
    Z_{Th, B} = \frac{R}{1 + 2u}
  \end{gathered}
\end{equation*}
\'E un partitore di tensione quindi
\begin{equation*}
  \begin{gathered}
    V_+ = \frac{Z_{Th, B} + R}{Z_{Th, A} + Z_{Th, B} + R + \frac{1}{sC}}V_{Th, A}
  \end{gathered}
\end{equation*}

\paragraph{Richiesta (c)}
Verificare che i limiti asintotici del guadagno (sia per $f \to 0$ che per $f \to \infty$) coincidano
con quanto osservato e qualitativamente dedotto al punto 2a.

\paragraph{Richiesta (d)}
Dall'espressione ottenuta dedurre i valori attesi per i parametri del filtro (frequenza
centrale, larghezza di banda e fattore di merito) e confrontarli con le misure di cui al
punto 1b (si consiglia di utilizzare la grandezza adimensionale $x = \omega RC = 2\pi fRC$
come variabile indipendente e di semplificare preliminarmente l'espressione della funzione
di trasferimento dividendo per il numeratore).

\paragraph{Richiesta (e)}
Osservare le propriet\'a di simmetria del plot di Bode e giustificarle alla luce
del fatto che $A(1/x) = A^*(x)$.
